\documentclass[a4paper,12pt]{article}
\usepackage[T2A]{fontenc}
\usepackage[utf8]{inputenc}
\usepackage[russian]{babel}
\usepackage{amsmath,amssymb,mathtools}
\usepackage{helvet}
\renewcommand{\familydefault}{\sfdefault}
\usepackage{icomma}
\usepackage[margin=19mm,top=21mm,bottom=22mm]{geometry}
\usepackage{microtype}
\usepackage{xcolor}
\usepackage{tikz}
\usetikzlibrary{arrows.meta,positioning,shapes.geometric,calc}
\usepackage{tabularx,booktabs}
\usepackage{enumitem}
\usepackage{fancyhdr}
\usepackage{setspace}
\usepackage{hyperref}
\usepackage{needspace}
\definecolor{accent}{HTML}{17696E}
\definecolor{deepaccent}{HTML}{12464A}
\definecolor{textgray}{HTML}{37434A}
\definecolor{lineblue}{HTML}{568DB2}
\hypersetup{colorlinks=true,linkcolor=deepaccent,urlcolor=accent}
\onehalfspacing
\setlength{\parindent}{0pt}
\setlength{\parskip}{4pt}
\setlist[itemize]{topsep=2pt,itemsep=2pt,leftmargin=6mm}
\setlist[enumerate]{topsep=2pt,itemsep=7pt,leftmargin=8mm}
\setcounter{secnumdepth}{0}
\renewcommand{\thesection}{}
\newcommand{\checkitem}{\(\square\)\enspace}
\newcommand{\important}[1]{\textcolor{deepaccent}{\textbf{#1}}}
\newcommand{\method}[1]{\medskip\noindent\textcolor{deepaccent}{\textbf{#1}}\quad}
\pagestyle{fancy}
\fancyhf{}
\fancyhead[L]{\footnotesize\color{textgray} ЕГЭ · Текстовые задачи и математические модели}
\fancyhead[R]{\footnotesize\color{textgray} 21.10.2026}
\fancyfoot[C]{\footnotesize\color{textgray} \thepage}
\renewcommand{\headrulewidth}{0.25pt}
\renewcommand{\footrulewidth}{0pt}
\setlength{\headheight}{15pt}
\newcommand{\sectiontitle}[1]{\needspace{6\baselineskip}\section*{\color{deepaccent}#1}}
\begin{document}

% Титульный лист
\thispagestyle{empty}
\begin{center}
\vspace*{28mm}
{\color{accent}\large МАТЕМАТИКА · ПОДГОТОВКА К ЕГЭ}\par
\vspace{21mm}
{\Huge\bfseries Чек-лист}\par
\vspace{8mm}
{\LARGE\color{deepaccent}\bfseries Прикладные и текстовые задачи}\par
\vspace{6mm}
{\large Физические модели · сплавы · высушивание}\par
\vspace{17mm}
\rule{0.65\textwidth}{0.5pt}\par
\vspace{10mm}
{\Large Ксения Клыкова}\par
\vspace{4mm}
{\large 21 октября 2026 года}\par
\vfill
{\normalsize\color{textgray} Лёвин Артём Александрович\\[3pt] эксклюзивно для Ксении Клыковой}\par
\vspace{12mm}
\begin{tikzpicture}[x=1cm,y=1cm]
\draw[accent!65,line width=1.2pt] (-5,0) .. controls (-2.8,0.65) and (-1.1,-0.42) .. (0.45,0.03)
  .. controls (2.1,0.47) and (3.2,-0.14) .. (5,0.30);
\end{tikzpicture}
\vspace*{7mm}
\end{center}
\clearpage

\sectiontitle{Маршрут повторения}
Отметьте пункт, когда сможете воспроизвести соответствующий приём без подсказки.
\begin{itemize}[label={}]
\item \checkitem По двум временам падения вычисляю изменение высоты, используя разность квадратов.
\item \checkitem Подставляю величины в физическую формулу, перевожу единицы в СИ, отбираю допустимое значение скорости.
\item \checkitem Связываю долю объёма и долю высоты цилиндрического бака и выбираю физически возможное время.
\item \checkitem Записываю баланс массы сплавов и массу чистого компонента через концентрацию.
\item \checkitem В задачах на высушивание сохраняю массу сухого вещества, а не массу воды.
\end{itemize}

\sectiontitle{1. Подъём воды в колодце}
\important{Исходная зависимость:} если камень падает без начальной скорости, то
\[
  h=\frac{gt^2}{2}=5t^2\qquad (g=10\ \text{м/с}^2).
\]
Время падения до дождя \(t_1=1{,}3\) с; после дождя оно уменьшилось на \(0{,}2\) с, то есть \(t_2=1{,}1\) с. Вода поднялась на разность пройденных камнем расстояний:
\[
\Delta h=h_1-h_2=5\bigl(t_1^2-t_2^2\bigr)
=5(t_1-t_2)(t_1+t_2)=5\cdot0{,}2\cdot2{,}4=\boxed{2{,}4\ \text{м}}.
\]
\begin{center}
\begin{tikzpicture}[>=Latex,scale=.95,every node/.style={font=\small}]
  \foreach \xx/\lab/\water/\height in {0/{До дождя}/0.35/{\(h_1\)},5/{После дождя}/0.9/{\(h_2\)}}{
    \draw[deepaccent,thick] (\xx,2.25) -- (\xx,0) -- (\xx+2.4,0) -- (\xx+2.4,2.25);
    \fill[lineblue!16] (\xx+.02,0.02) rectangle (\xx+2.38,\water);
    \draw[lineblue,very thick] (\xx,\water) -- (\xx+2.4,\water);
    \draw[accent,very thick,<-] (\xx+1.2,\water+0.05) -- (\xx+1.2,2.18);
    \node[anchor=west] at (\xx+1.25,1.6){\height};
    \node[above] at (\xx+1.2,2.35){\lab};
  }
  \draw[dashed,textgray] (0,2.25)--(7.4,2.25);
\end{tikzpicture}
\end{center}
\method{Быстрый приём.} В выражении \(t_1^2-t_2^2\) сначала используйте формулу разности квадратов: это уменьшает объём вычислений и риск арифметической ошибки.

\sectiontitle{2. Ведро на верёвке: граничное условие}
В верхней точке окружности давление воды на дно ведра описано формулой
\[
P=m\left(\frac{v^2}{L}-g\right),
\]
где \(m>0\) --- масса воды, \(v\) --- скорость, \(L\) --- длина верёвки. Чтобы вода не выливалась, в верхней точке должно выполняться \(P\ge0\). \important{Минимальная скорость} соответствует \(P=0\):
\[
\frac{v^2}{L}-g=0\quad\Longrightarrow\quad
v_{\min}=\sqrt{gL}.
\]
Если \(L=40\) см \(=0{,}4\) м и \(g=10\ \text{м/с}^2\), то
\[
v^2=10\cdot0{,}4=4,\qquad v_{\min}=\boxed{2\ \text{м/с}}.
\]
\method{Проверьте смысл.} Масса сокращается только потому, что \(m>0\); \(v\) --- модуль скорости, поэтому отрицательное значение не подходит.

\sectiontitle{3. Опорожнение цилиндрического бака}
После открытия крана высота столба воды меняется по закону
\[
H(t)=H_0-\sqrt{2gH_0}\,kt+\frac{g}{2}k^2t^2,
\]
где \(H_0\) --- начальная высота, \(k>0\) --- постоянный коэффициент, \(t\ge0\) --- время. \important{Для бака с постоянной площадью поперечного сечения} \(V=SH\). Следовательно, если осталась четверть первоначального объёма, то \(H=H_0/4\).

При \(H_0=20\) м, \(k=1/50\), \(g=10\ \text{м/с}^2\) получаем \(H=5\) м:
\begin{align*}
5&=20-\sqrt{2\cdot10\cdot20}\cdot\frac{t}{50}
       +\frac{10}{2}\cdot\frac{t^2}{50^2},\\
0&=15-\frac{2t}{5}+\frac{t^2}{500},\\
0&=7500-200t+t^2,\\
t^2-200t+7500&=0,\quad D=40000-30000=10000,\\
t_1&=50\ \text{с},\qquad t_2=150\ \text{с}.
\end{align*}
\needspace{13\baselineskip}\important{Почему берём 50 секунд?} Высота в данном законе может быть представлена как квадрат:
\[
H(t)=\left(\sqrt{H_0}-k\sqrt{\frac g2}\,t\right)^2.
\]
Физическая модель действует до опорожнения бака. Оно происходит при
\[
T_{\text{пуст}}=\frac{1}{k}\sqrt{\frac{2H_0}{g}}
=50\sqrt{\frac{40}{10}}=100\ \text{с}.
\]
Таким образом, \(0\le t\le100\) с. Значение \(150\) с выходит за пределы физического процесса, а \(50\) с допустимо. \(\boxed{t=50\ \text{с}}\)

\method{Универсальный способ.} Задайте нужную высоту \(H\), решите уравнение относительно \(t\), затем обязательно проверьте \(0\le t\le T_{\text{пуст}}\). Если \(H=\alpha H_0\), где \(0\le\alpha\le1\), то непосредственно получаем
\[
\boxed{\displaystyle t=\frac{1-\sqrt\alpha}{k}\sqrt{\frac{2H_0}{g}}}\qquad
\text{(физически допустимый момент)}.
\]

\sectiontitle{4. Сплавы: баланс двух масс}
В смеси сохраняются \important{общая масса} и \important{масса нужного компонента}. Концентрацию \(p\%\) записывайте долей \(p/100\), а массу вещества вычисляйте по формуле \(m_{\text{в-ва}}=c\cdot m_{\text{сплава}}\).

Два сплава содержат соответственно \(5\%\) и \(30\%\) никеля. Из них получили \(125\) кг сплава с \(25\%\) никеля. Пусть массы исходных сплавов равны \(x\) и \(y\) кг.
\begin{center}
\renewcommand{\arraystretch}{1.22}
\begin{tabularx}{\linewidth}{@{}Xccc@{}}
\toprule
Сплав & Масса, кг & Доля никеля & Никель, кг\\
\midrule
Первый & \(x\) & \(0{,}05\) & \(0{,}05x\)\\
Второй & \(y\) & \(0{,}30\) & \(0{,}30y\)\\
Смесь & \(125\) & \(0{,}25\) & \(0{,}25\cdot125\)\\
\bottomrule
\end{tabularx}
\end{center}
\[
\begin{cases}
x+y=125,\\
0{,}05x+0{,}30y=0{,}25\cdot125.
\end{cases}
\quad\Longrightarrow\quad
\boxed{x=25\ \text{кг},\ y=100\ \text{кг}}.
\]
\method{Вариант модели.} Если второй сплав тяжелее первого на \(3\) кг, обозначьте массы как \(x\) и \(x+3\); их сумму --- как \(2x+3\). Например, для долей меди \(0{,}10\), \(0{,}40\), \(0{,}30\) составьте
\[
0{,}10x+0{,}40(x+3)=0{,}30(2x+3).
\]
Отсюда \(x=3\), второй сплав весит \(6\) кг, итоговая масса \(\boxed{9\ \text{кг}}\).

\sectiontitle{5. Высушивание: что сохраняется?}
При высушивании \important{вода испаряется, а масса сухого вещества сохраняется}. Если виноград содержит \(90\%\) воды, то сухого вещества в нём \(10\%\). В изюме с \(5\%\) воды сухого вещества \(95\%\).

Для получения \(60\) кг изюма из \(x\) кг винограда:
\[
\underbrace{0{,}10x}_{\text{сухое вещество до}}
=\underbrace{0{,}95\cdot60}_{\text{сухое вещество после}},
\qquad x=\boxed{570\ \text{кг}}.
\]
\method{Контрольный вопрос.} Какое вещество меняет свою массу? Вода. Какое остаётся неизменным? Мякоть (сухое вещество). Именно его массу нужно приравнивать.

\sectiontitle{6. Четыре обязательные проверки}
\begin{itemize}
\item \checkitem \textbf{Размерности.} Сантиметры переведены в метры; время --- в секунды.
\item \checkitem \textbf{Проценты.} \(5\%=0{,}05\), \(90\%=0{,}90\), а не \(5\) и \(90\).
\item \checkitem \textbf{Преобразования.} Домножая \emph{уравнение} на ненулевой общий знаменатель, умножайте каждое слагаемое обеих частей.
\item \checkitem \textbf{Отбор.} Положительные длины и массы, неотрицательная скорость, время в допустимом физическом интервале.
\end{itemize}

\clearpage
\thispagestyle{fancy}
\begin{center}
{\Large\bfseries\color{deepaccent} Алгоритм: от условия к допустимому ответу}
\end{center}
\vspace{8mm}
\begin{center}
\begin{tikzpicture}[
>=Latex,
node distance=11mm,
actionbox/.style={rectangle,rounded corners=3pt,draw=accent,thick,align=center,text width=108mm,minimum height=15mm,inner sep=3mm,font=\small},
choice/.style={diamond,aspect=3.15,draw=accent,thick,align=center,text width=74mm,inner sep=1.6mm,font=\small},
terminal/.style={ellipse,draw=deepaccent,thick,align=center,text width=72mm,minimum height=12mm,font=\small},
arrow/.style={->,line width=.9pt,draw=deepaccent},
branch/.style={font=\small,fill=white,inner sep=2pt}]
\node[terminal] (start) {Прочитайте условие и определите искомое};
\node[actionbox,below=of start] (quantity) {Запишите данные с единицами измерения;\\назовите сохраняющуюся или связанную величину};
\node[actionbox,below=of quantity] (model) {Постройте математическую модель:\\формулу, таблицу, уравнение или систему};
\node[actionbox,below=of model] (sol) {Получите все алгебраические кандидаты;\\выберите очередной и проверьте вычисления};
\node[choice,below=18mm of sol] (valid) {Кандидат удовлетворяет\\ограничениям задачи?};
\node[actionbox,below=18mm of valid] (accept) {Оставьте допустимый корень;\\проверьте подстановкой и смыслом};
\node[terminal,below=of accept] (finish) {Запишите ответ с единицами измерения};
\node[actionbox,draw=textgray,text width=30mm,minimum height=18mm,right=13mm of valid] (reject) {Отбросьте корень;\\проверьте следующий};
\draw[arrow] (start)--(quantity);
\draw[arrow] (quantity)--(model);
\draw[arrow] (model)--(sol);
\draw[arrow] (sol)--(valid);
\draw[arrow] (valid)--node[branch,left]{да}(accept);
\draw[arrow] (accept)--(finish);
\draw[arrow] (valid)--node[branch,above]{нет}(reject);
\draw[arrow] (reject.north) |- ([yshift=-5mm]sol.east) -- (sol.east);
\end{tikzpicture}
\end{center}
\clearpage

\sectiontitle{Тренировочный блок · Путь А}
\textcolor{textgray}{10 задач. Выполняйте последовательно; решения и промежуточные вычисления оформляйте самостоятельно.}
\begin{enumerate}[label=\textbf{\arabic*.},wide=0pt]
\item До дождя камень падал в колодец \(1{,}5\) с. После дождя время падения стало на \(0{,}1\) с меньше. На сколько метров поднялась вода? Считать \(h=5t^2\).

\item Камень падал до поверхности воды \(1{,}8\) с, а после подъёма уровня воды --- \(1{,}5\) с. На сколько метров изменился уровень воды? Считать \(g=10\ \text{м/с}^2\).

\item Ведро с водой вращают в вертикальной плоскости на верёвке длиной \(90\) см. Какова наименьшая скорость в верхней точке, при которой вода не выливается? \(g=10\ \text{м/с}^2\).

\item Скорость ведра в верхней точке равна \(5\) м/с. Какова \emph{наибольшая} длина верёвки, при которой вода не выливается? \(g=10\ \text{м/с}^2\).

\item Высота воды в цилиндрическом баке изменяется по закону
\[
H(t)=H_0-\sqrt{2gH_0}\,kt+\frac g2 k^2t^2.
\]
Найдите, через сколько секунд останется четверть начального объёма, если \(H_0=20\) м, \(g=10\ \text{м/с}^2\), \(k=1/40\).

\item Для другого цилиндрического бака действует та же формула. Дано: \(H_0=45\) м, \(g=10\ \text{м/с}^2\), \(k=1/30\). Через сколько секунд останется \(1/9\) первоначального объёма? Отберите физически допустимое время.

\item Смешали два сплава с \(5\%\) и \(30\%\) никеля. Получили \(100\) кг сплава, содержащего \(20\%\) никеля. Определите массы обоих исходных сплавов.

\item Первый сплав содержит \(15\%\) меди, второй --- \(45\%\). Второй сплав на \(4\) кг тяжелее первого. После смешивания концентрация меди стала \(35\%\). Определите массу смеси.

\item В винограде \(86\%\) воды, а в изюме \(16\%\) воды. Сколько килограммов изюма получится из \(300\) кг винограда, если масса сухого вещества при высушивании не меняется?

\item В свежем винограде \(90\%\) воды, а в изюме \(5\%\) воды. Сколько килограммов винограда требуется для получения \(38\) кг изюма? 
\end{enumerate}

\clearpage
\sectiontitle{Ответы и самопроверка}
\renewcommand{\arraystretch}{1.43}
\begin{tabularx}{\linewidth}{@{}c>{\raggedright\arraybackslash}p{29mm}X@{}}
\toprule
№ & Ответ & Ключевой шаг для проверки\\
\midrule
1 & \(1{,}45\) м & \(5(1{,}5^2-1{,}4^2)\).\\
2 & \(4{,}95\) м & \(5(1{,}8^2-1{,}5^2)\).\\
3 & \(3\) м/с & \(v_{\min}=\sqrt{10\cdot0{,}9}\).\\
4 & \(2{,}5\) м & \(L\le v^2/g=25/10\).\\
5 & \(40\) с & Время опорожнения \(80\) с; для \(H_0/4\) проходит половина.\\
6 & \(60\) с & \(H=5\); корни \(60\) и \(120\) с, а опорожнение --- в \(90\) с.\\
7 & \(40\) кг и \(60\) кг & \(x+y=100\); \(0{,}05x+0{,}30y=20\).\\
8 & \(12\) кг & \(0{,}15x+0{,}45(x+4)=0{,}35(2x+4)\), \(x=4\).\\
9 & \(50\) кг & \(0{,}14\cdot300=0{,}84m\).\\
10 & \(361\) кг & \(0{,}10x=0{,}95\cdot38\).\\
\bottomrule
\end{tabularx}

\bigskip
\important{Итоговая самопроверка.}
\begin{itemize}
\item Если уравнение дало несколько значений, перечислите все кандидаты и укажите, какие именно условия отсекли лишние.
\item Для смесей пересчитайте суммарную массу чистого вещества двумя способами.
\item Для высушивания сравните массу сухого вещества до и после --- она должна совпасть.
\end{itemize}
\end{document}
